## TEMP holding note — RQ1 physics-weight ablation (2026-08-11)

**Fit ran on cluster, evaluated locally**: 2026-08-11
(`jobs/rq123_methodology/rq1_physicsablation_fit.sh`, 32 jobs — see that file's own header for
the full design rationale). All 32/32 completed, zero `FAILED` entries. Two tiers: a cheap
single-split exploratory pass across `physics_weight`/`trajectory_weight` ∈ {0, 1, 2}, and a
rigorous 5-fold `spatial_block_kfold` pass at `physics_weight=0` only (comparable directly to the
existing `physics_weight=1` kfold numbers already in `TEMP_rq1_sweep_results_2026-08-11.tex`).
DNN is not included — it has no physics loss or `y_max`/`k` sub-network at all, so it isn't
something this ablation can vary; its own numbers from the main sweep are the implicit
zero-architecture reference point.

### How these numbers were generated

- Single-split tier: `run_pinn_env_terrain.py`/`run_pinn_env_terrain_k.py`,
  `--split-type spatial_block`, `--physics-weight`/`--trajectory-weight` set together to 0.0/1.0/2.0,
  Set3 (`nested_set3_gated_terrain_wind_vif`), both cohorts. Evaluated with each model's own
  `evaluate_*` script.
- 5-fold tier: same models, `--split-type spatial_block_kfold`, `physics_weight=trajectory_weight=0.0`
  only (the `w=1` kfold numbers already existed from the main sweep). Pooled with
  `models.common.kfold_summary`, same convention as the main sweep (pooled R2 over the whole
  population + bootstrap CI + per-fold mean/SD).

### Results — single-split exploratory tier (weights 0/1/2)

| Model | Cohort | w=0 R2 | w=1 R2 | w=2 R2 |
|---|---|---:|---:|---:|
| PINN | 4survey | 0.630 | 0.584 | 0.561 |
| PINN_k | 4survey | 0.632 | 0.587 | 0.574 |
| PINN | 6survey | 0.747 | 0.731 | 0.723 |
| PINN_k | 6survey | 0.747 | 0.725 | 0.720 |

### Results — 5-fold rigorous tier (w=0 vs. the existing w=1 kfold numbers)

| Model | Cohort | w=0 pooled R2 [95% CI] | w=1 pooled R2 [95% CI] | Δ (w0 − w1) |
|---|---|---|---|---:|
| PINN | 4survey | 0.6295 [0.600, 0.653] | 0.5775 [0.545, 0.607] | **+0.052** |
| PINN | 6survey | 0.7199 [0.685, 0.744] | 0.7097 [0.671, 0.736] | +0.010 |
| PINN_k | 4survey | 0.6288 [0.599, 0.653] | 0.5788 [0.547, 0.608] | **+0.050** |
| PINN_k | 6survey | 0.7187 [0.686, 0.741] | 0.7083 [0.671, 0.735] | +0.010 |

### Headline finding

**Monotonic, no exceptions: R2 decreases as physics_weight increases, across both tiers, both
models, both cohorts.** The single-split exploratory tier (0/1/2) shows a smooth monotonic
decline; the rigorous 5-fold tier confirms w=0 beats w=1 by a real margin on 4survey (+0.05,
comfortably outside both configurations' 95% CIs — not noise) and a smaller but still positive
margin on 6survey (+0.01, within the wider 6survey CIs — real direction, not a strong claim on
6survey specifically). **This directly answers "why doesn't PINN improve over DNN"**: it's not
that the `y_max`/`k` sub-network architecture is useless — it's that the physics LOSS specifically
is actively counterproductive at this project's data scale, not just neutral. DNN's own numbers
(no architecture, no physics loss at all) still beat every PINN configuration including w=0 on
4survey (DNN Set3 = 0.659 vs. PINN w=0 = 0.630) — so removing the physics loss closes most, but
not all, of the DNN-PINN gap on 4survey; on 6survey PINN w=0 (0.720/0.719) essentially closes the
gap with DNN entirely (DNN Set3 6survey = 0.707).

**Secondary finding, PINN_k's y_max/k correlation gets WORSE without the physics constraint** —
at w=1 (main sweep): 4survey=-0.449, 6survey=-0.057. At w=0: 4survey=**-0.935** (near-perfect
anti-correlation, the two learned parameters become almost interchangeable), 6survey=**+0.430**
(flips sign entirely). This is a real accuracy-vs-identifiability tradeoff, not a clean win for
removing the physics constraint: w=0 gives better raw R2 but a substantially less interpretable
model, since y_max and k become much harder to attribute independently once nothing is pulling
them toward physically-consistent, separable values.
